% ============================================================================
%  MANUEL UTILISATEUR — EcoleGest v5.0
%  Application de Gestion Scolaire Intégrée
%  Melen, Yaoundé — Cameroun
% ============================================================================
\documentclass[12pt, a4paper, twoside, openright]{report}

% ─── ENCODAGE & LANGUE ──────────────────────────────────────────────────────
\usepackage[utf8]{inputenc}
\usepackage[T1]{fontenc}
\usepackage[french]{babel}

% ─── MISE EN PAGE ────────────────────────────────────────────────────────────
\usepackage[
  top=2.5cm, bottom=2.5cm,
  left=2.8cm, right=2.2cm,
  headheight=15pt
]{geometry}
\usepackage{setspace}
\onehalfspacing

% ─── POLICES ─────────────────────────────────────────────────────────────────
\usepackage{lmodern}
\usepackage{microtype}

% ─── GRAPHIQUES & COULEURS ───────────────────────────────────────────────────
\usepackage{graphicx}
\usepackage[dvipsnames, svgnames, x11names, table]{xcolor}
\usepackage{tikz}
\usetikzlibrary{shadows, calc, positioning, arrows.meta}

% ─── TABLEAUX ────────────────────────────────────────────────────────────────
\usepackage{booktabs}
\usepackage{tabularx}
\usepackage{multirow}
\usepackage{array}
\usepackage{colortbl}

% ─── LISTES ──────────────────────────────────────────────────────────────────
\usepackage{enumitem}

% ─── EN-TÊTES & PIEDS DE PAGE ───────────────────────────────────────────────
\usepackage{fancyhdr}
\pagestyle{fancy}
\fancyhf{}
\fancyhead[LE]{\small\textit{\leftmark}}
\fancyhead[RO]{\small\textit{\rightmark}}
\fancyfoot[C]{\thepage}
\fancyfoot[LE,RO]{\small\textcolor{gray}{EcoleGest — Manuel Utilisateur v5.0}}
\renewcommand{\headrulewidth}{0.4pt}
\renewcommand{\footrulewidth}{0.2pt}

% ─── HYPERLIENS ──────────────────────────────────────────────────────────────
\usepackage[
  colorlinks=true,
  linkcolor=NavyBlue,
  urlcolor=RoyalBlue,
  citecolor=ForestGreen,
  bookmarks=true,
  bookmarksopen=true,
  pdfauthor={Direction Technique — EcoleGest},
  pdftitle={Manuel Utilisateur EcoleGest v5.0},
  pdfsubject={Guide complet d'utilisation},
  pdfkeywords={école, gestion, scolaire, EcoleGest, Cameroun}
]{hyperref}
\usepackage{bookmark}

% ─── BOÎTES COLORÉES (tips, warnings, notes…) ───────────────────────────────
\usepackage[most]{tcolorbox}

% ─── ICÔNES ──────────────────────────────────────────────────────────────────
\usepackage{fontawesome5}

% ─── LISTINGS CODE ───────────────────────────────────────────────────────────
\usepackage{listings}

% ============================================================================
%  COULEURS PERSONNALISÉES (palette EcoleGest)
% ============================================================================
\definecolor{brandBlue}{HTML}{0062FF}
\definecolor{brandOrange}{HTML}{FFA000}
\definecolor{bgLight}{HTML}{F8FAFC}
\definecolor{bgDark}{HTML}{0F172A}
\definecolor{success}{HTML}{059669}
\definecolor{danger}{HTML}{DC2626}
\definecolor{warning}{HTML}{D97706}
\definecolor{info}{HTML}{0EA5E9}
\definecolor{grayMedium}{HTML}{64748B}
\definecolor{grayLight}{HTML}{E2E8F0}

% ============================================================================
%  BOÎTES PERSONNALISÉES
% ============================================================================

% ── Astuce ───────────────────────────────────────────────────────────────────
\newtcolorbox{astuce}[1][]{
  enhanced, breakable,
  colback=success!5, colframe=success!80!black,
  fonttitle=\bfseries\color{white},
  title={\faLightbulb\quad Astuce},
  attach boxed title to top left={yshift=-2mm, xshift=4mm},
  boxed title style={colback=success!80!black, rounded corners},
  rounded corners, arc=3mm,
  left=4mm, right=4mm, top=6mm, bottom=4mm,
  #1
}

% ── Attention ────────────────────────────────────────────────────────────────
\newtcolorbox{attention}[1][]{
  enhanced, breakable,
  colback=danger!5, colframe=danger!80!black,
  fonttitle=\bfseries\color{white},
  title={\faExclamationTriangle\quad Attention},
  attach boxed title to top left={yshift=-2mm, xshift=4mm},
  boxed title style={colback=danger!80!black, rounded corners},
  rounded corners, arc=3mm,
  left=4mm, right=4mm, top=6mm, bottom=4mm,
  #1
}

% ── Note ─────────────────────────────────────────────────────────────────────
\newtcolorbox{note}[1][]{
  enhanced, breakable,
  colback=info!5, colframe=info!80!black,
  fonttitle=\bfseries\color{white},
  title={\faInfoCircle\quad Note},
  attach boxed title to top left={yshift=-2mm, xshift=4mm},
  boxed title style={colback=info!80!black, rounded corners},
  rounded corners, arc=3mm,
  left=4mm, right=4mm, top=6mm, bottom=4mm,
  #1
}

% ── Étape (pour les procédures pas-à-pas) ────────────────────────────────────
\newtcolorbox{etape}[1][]{
  enhanced, breakable,
  colback=brandBlue!3, colframe=brandBlue!60!black,
  fonttitle=\bfseries\color{white},
  title={#1},
  attach boxed title to top left={yshift=-2mm, xshift=4mm},
  boxed title style={colback=brandBlue!60!black, rounded corners},
  rounded corners, arc=2mm,
  left=4mm, right=4mm, top=6mm, bottom=4mm,
}

% ============================================================================
%  COMMANDE POUR LES CAPTURES D'ÉCRAN
% ============================================================================
% Utilisation: \screenshot{nom_fichier}{Légende descriptive}
% Le fichier doit se trouver dans le dossier screenshots/
% Si l'image n'existe pas encore, une boîte placeholder s'affiche.
\newcommand{\screenshot}[2]{%
  \begin{figure}[htbp]
    \centering
    \IfFileExists{screenshots/#1}{%
      \includegraphics[width=0.92\textwidth, keepaspectratio]{screenshots/#1}%
    }{%
      \begin{tcolorbox}[
        enhanced, width=0.92\textwidth,
        height=7cm, valign=center,
        colback=bgLight, colframe=grayLight,
        fontupper=\centering\color{grayMedium}\large,
        rounded corners, arc=4mm,
        shadow={2mm}{-2mm}{0mm}{black!15}
      ]
        \faCamera\quad\textbf{Capture d'écran}\\[6pt]
        \normalsize\texttt{screenshots/\detokenize{#1}}\\[4pt]
        \small\textit{#2}
      \end{tcolorbox}%
    }%
    \caption{#2}
    \label{fig:#1}
  \end{figure}
}

% Variante pleine page
\newcommand{\screenshotfull}[2]{%
  \begin{figure}[p]
    \centering
    \IfFileExists{screenshots/#1}{%
      \includegraphics[width=\textwidth, keepaspectratio]{screenshots/#1}%
    }{%
      \begin{tcolorbox}[
        enhanced, width=\textwidth,
        height=12cm, valign=center,
        colback=bgLight, colframe=grayLight,
        fontupper=\centering\color{grayMedium}\Large,
        rounded corners, arc=4mm,
        shadow={2mm}{-2mm}{0mm}{black!15}
      ]
        \faDesktop\quad\textbf{Capture d'écran — Pleine page}\\[8pt]
        \large\texttt{screenshots/\detokenize{#1}}\\[6pt]
        \normalsize\textit{#2}
      \end{tcolorbox}%
    }%
    \caption{#2}
    \label{fig:#1}
  \end{figure}
}

% ============================================================================
%  COMMANDES UTILITAIRES
% ============================================================================
\newcommand{\menu}[1]{\textbf{\textcolor{brandBlue}{\faChevronRight~#1}}}
\newcommand{\bouton}[1]{\tikz[baseline=(X.base)]{\node[fill=brandBlue, text=white, rounded corners=3pt, inner sep=3pt, font=\small\bfseries] (X) {#1};}}
\newcommand{\boutonvert}[1]{\tikz[baseline=(X.base)]{\node[fill=success, text=white, rounded corners=3pt, inner sep=3pt, font=\small\bfseries] (X) {#1};}}
\newcommand{\boutonrouge}[1]{\tikz[baseline=(X.base)]{\node[fill=danger, text=white, rounded corners=3pt, inner sep=3pt, font=\small\bfseries] (X) {#1};}}
\newcommand{\champ}[1]{\texttt{\textcolor{bgDark}{#1}}}
\newcommand{\role}[1]{\textsc{\textcolor{brandOrange}{#1}}}

% Numérotation des chapitres plus propre
\renewcommand{\thechapter}{\arabic{chapter}}

% ============================================================================
%                         DÉBUT DU DOCUMENT
% ============================================================================
\begin{document}

% ============================================================================
%  PAGE DE COUVERTURE
% ============================================================================
\begin{titlepage}
  \begin{tikzpicture}[remember picture, overlay]
    % Fond dégradé
    \fill[brandBlue] (current page.south west) rectangle (current page.north east);
    \fill[brandOrange, opacity=0.15] (current page.south west) rectangle ([yshift=8cm]current page.east);
    
    % Motifs décoratifs
    \foreach \i in {1,...,20}{
      \fill[white, opacity=0.03] (rand*21, rand*29.7) circle (rand*2+0.5);
    }
    
    % Bande latérale
    \fill[brandOrange] ([xshift=-2mm]current page.north east) rectangle ([xshift=2mm]current page.south east);
    
    % Logo / Icône
    \node[circle, fill=white, minimum size=3cm, inner sep=0pt,
          drop shadow={shadow xshift=2mm, shadow yshift=-2mm, color=black!30}] 
      at ([yshift=-5cm]current page.north) {
        \IfFileExists{screenshots/logo.png}{%
          \includegraphics[width=2cm]{screenshots/logo.png}%
        }{%
          {\Huge\textcolor{brandBlue}{\faGraduationCap}}%
        }
      };
  \end{tikzpicture}
  
  \vspace*{6.5cm}
  \begin{center}
    {\fontsize{42}{48}\selectfont\bfseries\textcolor{white}{EcoleGest}}\\[8pt]
    {\Large\textcolor{white!85}{Application de Gestion Scolaire Intégrée}}\\[30pt]
    
    \rule{8cm}{1.5pt}\\[20pt]
    
    {\fontsize{22}{26}\selectfont\textcolor{brandOrange}{\textbf{MANUEL UTILISATEUR}}}\\[8pt]
    {\large\textcolor{white!80}{Guide Complet d'Utilisation — Version 5.0}}\\[40pt]
    
    \begin{tabular}{rl}
      \textcolor{white!70}{\faCalendarAlt} & \textcolor{white}{Juillet 2026}\\[4pt]
      \textcolor{white!70}{\faMapMarkerAlt} & \textcolor{white}{Melen, Yaoundé — Cameroun}\\[4pt]
      \textcolor{white!70}{\faFileAlt} & \textcolor{white}{Document confidentiel}\\
    \end{tabular}
  \end{center}
  
  \vfill
  \begin{center}
    \textcolor{white!50}{\small\textcopyright{} 2026 EcoleGest. Tous droits réservés.}
  \end{center}
\end{titlepage}

% ============================================================================
%  PAGE DE GARDE INTÉRIEURE
% ============================================================================
\thispagestyle{empty}
\vspace*{3cm}
\begin{center}
  {\LARGE\textbf{EcoleGest — Manuel Utilisateur}}\\[10pt]
  {\large Version 5.0 — Juillet 2026}\\[30pt]
  
  \begin{tabularx}{0.75\textwidth}{lX}
    \toprule
    \textbf{Propriété} & \textbf{Valeur} \\
    \midrule
    Titre & Manuel Utilisateur EcoleGest \\
    Version du document & 5.0 \\
    Version logiciel & 5.0 \\
    Date de publication & Juillet 2026 \\
    Classification & Confidentiel — Usage interne \\
    Auteur & Direction Technique \\
    Établissement & Complexe Scolaire EcoleGest \\
    Localisation & Melen, Yaoundé, Cameroun \\
    Contact support & support@ecolegest.cm \\
    \bottomrule
  \end{tabularx}
\end{center}

\vfill

\begin{tcolorbox}[colback=info!5, colframe=info, rounded corners, arc=3mm]
  \small
  \textbf{Historique des révisions}\\[6pt]
  \begin{tabularx}{\textwidth}{llXl}
    \textbf{Version} & \textbf{Date} & \textbf{Modifications} & \textbf{Auteur} \\
    \hline
    1.0 & Janvier 2026 & Création initiale & Équipe technique \\
    2.0 & Mars 2026 & Ajout modules enseignant/parent & Équipe technique \\
    3.0 & Mai 2026 & Module paiements et bulletins & Équipe technique \\
    4.0 & Juin 2026 & Messagerie et discipline & Équipe technique \\
    5.0 & Juillet 2026 & Version complète avec notifications & Équipe technique \\
  \end{tabularx}
\end{tcolorbox}

\clearpage

% ============================================================================
%  TABLE DES MATIÈRES
% ============================================================================
\tableofcontents
\clearpage

\listoffigures
\clearpage

% ============================================================================
%  CHAPITRE 1 : PRÉSENTATION GÉNÉRALE
% ============================================================================
\chapter{Présentation Générale}
\label{chap:presentation}

\section{Qu'est-ce que EcoleGest~?}

\textbf{EcoleGest} est une application web de gestion scolaire intégrée, conçue pour simplifier et automatiser l'ensemble des processus administratifs, pédagogiques et financiers d'un établissement scolaire.

L'application est accessible depuis n'importe quel navigateur web (Google Chrome, Mozilla Firefox, Microsoft Edge, Safari) sur ordinateur, tablette ou téléphone.

\begin{note}
  Aucune installation n'est nécessaire sur votre appareil. Il suffit d'ouvrir votre navigateur et de saisir l'adresse de l'application.
\end{note}

\section{Les acteurs de l'application}

EcoleGest est conçue pour \textbf{cinq types d'utilisateurs} (appelés « acteurs »). Chaque acteur dispose de son propre espace avec des fonctionnalités adaptées à son rôle.

\begin{table}[htbp]
  \centering
  \caption{Les acteurs de EcoleGest et leurs responsabilités}
  \label{tab:acteurs}
  \rowcolors{2}{bgLight}{white}
  \begin{tabularx}{\textwidth}{>{\bfseries}l X l}
    \toprule
    \textbf{Acteur} & \textbf{Responsabilités principales} & \textbf{Icône} \\
    \midrule
    Administrateur & 
      Gestion complète de l'école : élèves, enseignants, classes, cours, emploi du temps, bulletins, discipline, bus scolaire, personnel & 
      \faUserShield \\
    Directeur & 
      Supervision académique et financière, consultation des rapports, suivi des paiements, validation des bulletins & 
      \faUserTie \\
    Intendant & 
      Gestion financière : paiements, tranches de scolarité, suivi des impayés, inscriptions & 
      \faMoneyBillWave \\
    Enseignant & 
      Gestion pédagogique : emploi du temps personnel, saisie des notes, évaluations, cahier d'appel, communication & 
      \faChalkboardTeacher \\
    Parent & 
      Suivi de l'enfant : résultats scolaires, paiements, bulletins, communication avec l'école & 
      \faUsers \\
    \bottomrule
  \end{tabularx}
\end{table}

\section{Schéma général de l'application}

La figure ci-dessous présente une vue d'ensemble des modules de l'application et les acteurs qui y ont accès.

\screenshot{schema_general.png}{Vue d'ensemble des modules EcoleGest et des acteurs}

\section{Configuration requise}

\begin{table}[htbp]
  \centering
  \caption{Configuration minimale requise}
  \begin{tabularx}{\textwidth}{lX}
    \toprule
    \textbf{Élément} & \textbf{Exigence} \\
    \midrule
    Navigateur web & Google Chrome 90+, Firefox 88+, Edge 90+, Safari 14+ \\
    Connexion internet & Minimum 1 Mbps (recommandé : 5 Mbps) \\
    Écran & Résolution minimale 1024×768 pixels \\
    Système & Windows 10+, macOS 11+, Android 10+, iOS 14+ \\
    \bottomrule
  \end{tabularx}
\end{table}


% ============================================================================
%  CHAPITRE 2 : PRISE EN MAIN — CONNEXION
% ============================================================================
\chapter{Prise en Main}
\label{chap:prise_en_main}

\section{Accéder à l'application}

\begin{etape}{\faStepForward\quad Étape 1 — Ouvrir le navigateur}
  Ouvrez votre navigateur web préféré (Google Chrome est recommandé).\\[4pt]
  Dans la \textbf{barre d'adresse} (en haut du navigateur), tapez l'adresse de l'application :
  
  \begin{center}
    \Large\texttt{http://ecolegest.cm}
  \end{center}
  
  Puis appuyez sur la touche \textbf{Entrée} de votre clavier.
\end{etape}

\screenshot{landing_page.png}{Page d'accueil de EcoleGest — Page publique visible par tous}

\begin{note}
  La page d'accueil présente l'établissement, ses statistiques (nombre d'élèves, enseignants, classes), la localisation à Melen (Yaoundé) et les valeurs de rigueur et discipline.
\end{note}

\section{Se connecter à son compte}

Chaque utilisateur possède un \textbf{nom d'utilisateur} (identifiant) et un \textbf{mot de passe} qui lui sont attribués par l'administrateur de l'école.

\begin{etape}{\faStepForward\quad Étape 1 — Accéder à la page de connexion}
  Sur la page d'accueil, cliquez sur le bouton \bouton{Se Connecter} en haut à droite de l'écran.
\end{etape}

\screenshot{login_page.png}{Page de connexion — Saisie des identifiants}

\begin{etape}{\faStepForward\quad Étape 2 — Saisir ses identifiants}
  \begin{enumerate}[label=\textbf{\arabic*.}]
    \item Dans le champ \champ{Nom d'utilisateur}, tapez votre identifiant
    \item Dans le champ \champ{Mot de passe}, tapez votre mot de passe
    \item Cliquez sur le bouton \bouton{Connexion}
  \end{enumerate}
\end{etape}

\begin{attention}
  \begin{itemize}
    \item Votre mot de passe est \textbf{sensible aux majuscules/minuscules}. Vérifiez que la touche « Verr. Maj » (Caps Lock) de votre clavier est bien désactivée.
    \item \textbf{Ne partagez jamais} votre mot de passe avec une autre personne.
    \item Après 5 tentatives échouées, votre compte peut être temporairement bloqué. Contactez l'administrateur.
  \end{itemize}
\end{attention}

\begin{etape}{\faStepForward\quad Étape 3 — Vous êtes connecté !}
  Après une connexion réussie, vous serez automatiquement redirigé vers votre \textbf{tableau de bord} correspondant à votre rôle.
\end{etape}

\screenshot{login_success.png}{Redirection automatique vers le tableau de bord après connexion}

\section{Se déconnecter}

\begin{etape}{\faStepForward\quad Déconnexion}
  \begin{enumerate}[label=\textbf{\arabic*.}]
    \item Regardez le \textbf{menu latéral gauche} (la barre bleue et orange sur le côté)
    \item En bas de ce menu, vous verrez votre nom et un bouton \boutonrouge{Déconnexion}
    \item Cliquez sur \boutonrouge{Déconnexion}
    \item Vous serez redirigé vers la page de connexion
  \end{enumerate}
\end{etape}

\screenshot{sidebar_deconnexion.png}{Bouton de déconnexion dans le menu latéral}

\begin{astuce}
  Pensez toujours à vous déconnecter lorsque vous quittez votre poste, surtout si vous utilisez un ordinateur partagé.
\end{astuce}

\section{Comprendre l'interface}

L'interface de EcoleGest se compose de trois zones principales :

\begin{enumerate}[label=\textbf{\Alph*.}]
  \item \textbf{Le menu latéral (Sidebar)} — À gauche, contient tous les liens de navigation vers les différentes sections. Vous pouvez le réduire en cliquant sur la flèche.
  \item \textbf{La zone de contenu} — Au centre, affiche la page actuellement sélectionnée.
  \item \textbf{Les notifications} — Le badge rouge à côté de « Messages » indique le nombre de messages non lus.
\end{enumerate}

\screenshot{interface_zones.png}{Les trois zones de l'interface EcoleGest}


% ============================================================================
%  CHAPITRE 3 : ESPACE ADMINISTRATEUR
% ============================================================================
\chapter{Espace Administrateur}
\label{chap:admin}

L'administrateur est le \textbf{super-utilisateur} de l'application. Il a accès à toutes les fonctionnalités et est responsable de la configuration et de la gestion quotidienne de l'école.

\section{Tableau de bord}

Le tableau de bord administrateur affiche en un coup d'œil les indicateurs clés de l'établissement.

\screenshot{admin_dashboard.png}{Tableau de bord Administrateur — Vue d'ensemble}

\subsection*{Indicateurs affichés}
\begin{itemize}[label=\textcolor{brandBlue}{\faChartBar}]
  \item Nombre total d'élèves inscrits
  \item Nombre d'enseignants actifs
  \item Nombre de classes et salles
  \item Montant total des paiements reçus
  \item Graphiques de fréquentation et performance
\end{itemize}

% ─────────────────────────────────────────────────────────────────────────────
\section{Gestion des Élèves}
\label{sec:admin_eleves}

\subsection{Consulter la liste des élèves}

\begin{etape}{\faStepForward\quad Accéder à la liste}
  \begin{enumerate}[label=\textbf{\arabic*.}]
    \item Dans le menu latéral, cliquez sur \menu{Élèves}
    \item La liste de tous les élèves inscrits s'affiche dans un tableau
    \item Utilisez la \textbf{barre de recherche} en haut pour filtrer par nom ou matricule
    \item Utilisez le filtre \textbf{Classe} pour voir les élèves d'une classe spécifique
  \end{enumerate}
\end{etape}

\screenshot{admin_eleves_liste.png}{Liste des élèves — Vue tableau avec recherche et filtres}

\subsection{Ajouter un nouvel élève}

\begin{etape}{\faStepForward\quad Procédure d'ajout}
  \begin{enumerate}[label=\textbf{\arabic*.}]
    \item Cliquez sur le bouton \boutonvert{+ Ajouter un élève} en haut à droite
    \item Un formulaire apparaît — remplissez tous les champs marqués d'un astérisque (\textcolor{danger}{*}) :
      \begin{itemize}
        \item \champ{Nom} — Le nom de famille de l'élève
        \item \champ{Prénom} — Le ou les prénoms de l'élève
        \item \champ{Date de naissance} — Au format JJ/MM/AAAA
        \item \champ{Lieu de naissance} — La ville de naissance
        \item \champ{Sexe} — Masculin ou Féminin
        \item \champ{Classe} — Sélectionnez la classe dans la liste déroulante
      \end{itemize}
    \item Cliquez sur \bouton{Enregistrer}
    \item Un message de confirmation apparaît : « Élève ajouté avec succès »
  \end{enumerate}
\end{etape}

\screenshot{admin_eleves_ajout.png}{Formulaire d'ajout d'un nouvel élève}

\subsection{Modifier les informations d'un élève}

\begin{etape}{\faStepForward\quad Modification}
  \begin{enumerate}[label=\textbf{\arabic*.}]
    \item Dans la liste des élèves, trouvez l'élève concerné
    \item Cliquez sur l'icône \textcolor{brandBlue}{\faPencilAlt} (crayon) dans la colonne « Actions »
    \item Modifiez les champs souhaités
    \item Cliquez sur \bouton{Enregistrer}
  \end{enumerate}
\end{etape}

\subsection{Supprimer un élève}

\begin{attention}
  La suppression d'un élève est une action \textbf{irréversible}. Elle supprime également l'ensemble de ses évaluations, bulletins et paiements.
\end{attention}

\begin{etape}{\faStepForward\quad Suppression}
  \begin{enumerate}[label=\textbf{\arabic*.}]
    \item Dans la liste des élèves, trouvez l'élève concerné
    \item Cliquez sur l'icône \textcolor{danger}{\faTrash} (poubelle) dans la colonne « Actions »
    \item Une fenêtre de confirmation apparaît : « Êtes-vous sûr de vouloir supprimer cet élève~? »
    \item Cliquez sur \boutonrouge{Confirmer} pour supprimer, ou \bouton{Annuler} pour revenir
  \end{enumerate}
\end{etape}

% ─────────────────────────────────────────────────────────────────────────────
\section{Gestion des Enseignants}
\label{sec:admin_enseignants}

\subsection{Consulter la liste des enseignants}

\begin{etape}{\faStepForward\quad Accéder à la liste}
  \begin{enumerate}[label=\textbf{\arabic*.}]
    \item Dans le menu latéral, cliquez sur \menu{Enseignants}
    \item La liste affiche : nom, prénom, matière enseignée, nombre de classes, statut (actif/inactif)
  \end{enumerate}
\end{etape}

\screenshot{admin_enseignants_liste.png}{Liste des enseignants avec statistiques}

\subsection{Ajouter un enseignant}

\begin{etape}{\faStepForward\quad Procédure d'ajout}
  \begin{enumerate}[label=\textbf{\arabic*.}]
    \item Cliquez sur \boutonvert{+ Ajouter un enseignant}
    \item Remplissez le formulaire :
      \begin{itemize}
        \item \champ{Nom} et \champ{Prénom}
        \item \champ{Téléphone} — Numéro de téléphone mobile
        \item \champ{Nom d'utilisateur} — Sera utilisé pour la connexion (généré automatiquement si vide)
        \item \champ{Mot de passe} — Par défaut : \texttt{1234}
        \item \champ{Matière} — Sélectionnez le cours principal dans la liste
      \end{itemize}
    \item Cliquez sur \bouton{Enregistrer}
    \item \textbf{Important} : Communiquez les identifiants de connexion à l'enseignant
  \end{enumerate}
\end{etape}

\screenshot{admin_enseignants_ajout.png}{Formulaire d'ajout d'un enseignant}

\begin{astuce}
  Le nom d'utilisateur est généré automatiquement au format \texttt{prenom.nom}. Exemple : pour « Jean Dupont », le nom d'utilisateur sera \texttt{jean.dupont}.
\end{astuce}

\subsection{Affecter un enseignant à une classe/salle}

\begin{etape}{\faStepForward\quad Affectation}
  \begin{enumerate}[label=\textbf{\arabic*.}]
    \item Cliquez sur le nom de l'enseignant pour ouvrir sa fiche détaillée
    \item Cliquez sur \bouton{Affecter}
    \item Sélectionnez la \textbf{salle} et/ou le \textbf{cours} à attribuer
    \item Cliquez sur \bouton{Confirmer}
  \end{enumerate}
\end{etape}

\screenshot{admin_enseignants_affectation.png}{Affectation d'un enseignant à une classe}

% ─────────────────────────────────────────────────────────────────────────────
\section{Gestion des Classes et Salles}
\label{sec:admin_classes}

\subsection{Classes}

\begin{etape}{\faStepForward\quad Gérer les classes}
  \begin{enumerate}[label=\textbf{\arabic*.}]
    \item Dans le menu, cliquez sur \menu{Classes \& Salles}
    \item L'onglet \textbf{Classes} affiche toutes les classes avec le nombre d'élèves et le cycle
    \item Pour ajouter une classe, cliquez sur \boutonvert{+ Nouvelle classe}
    \item Renseignez le \champ{Libellé} (ex: « SIL », « CP », « CE1 ») et le \champ{Cycle}
    \item Cliquez sur \bouton{Enregistrer}
  \end{enumerate}
\end{etape}

\screenshot{admin_classes.png}{Gestion des classes — Liste et création}

\subsection{Salles}

\begin{etape}{\faStepForward\quad Gérer les salles}
  \begin{enumerate}[label=\textbf{\arabic*.}]
    \item Dans le menu, cliquez sur \menu{Salles}
    \item La liste des salles s'affiche avec la capacité et la classe associée
    \item Pour ajouter une salle, cliquez sur \boutonvert{+ Nouvelle salle}
    \item Renseignez le \champ{Nom de la salle}, la \champ{Capacité} et la \champ{Classe} associée
    \item Cliquez sur \bouton{Enregistrer}
  \end{enumerate}
\end{etape}

\screenshot{admin_salles.png}{Gestion des salles de classe}

% ─────────────────────────────────────────────────────────────────────────────
\section{Inscriptions}
\label{sec:admin_inscriptions}

\begin{etape}{\faStepForward\quad Inscrire un élève}
  \begin{enumerate}[label=\textbf{\arabic*.}]
    \item Dans le menu, cliquez sur \menu{Inscriptions}
    \item Cliquez sur \boutonvert{+ Nouvelle inscription}
    \item Sélectionnez l'\textbf{élève} dans la liste déroulante (ou créez-en un nouveau)
    \item Sélectionnez la \textbf{classe} et la \textbf{salle}
    \item Sélectionnez l'\textbf{année académique}
    \item Cliquez sur \bouton{Inscrire}
  \end{enumerate}
\end{etape}

\screenshot{admin_inscriptions.png}{Module d'inscription — Formulaire et liste}

% ─────────────────────────────────────────────────────────────────────────────
\section{Cours et Évaluations}
\label{sec:admin_cours}

\subsection{Gérer les cours (matières)}

\begin{etape}{\faStepForward\quad Ajouter un cours}
  \begin{enumerate}[label=\textbf{\arabic*.}]
    \item Dans le menu, cliquez sur \menu{Cours \& Évaluations}
    \item L'onglet \textbf{Cours} affiche les matières existantes
    \item Cliquez sur \boutonvert{+ Nouveau cours}
    \item Remplissez :
      \begin{itemize}
        \item \champ{Libellé} — Nom de la matière (ex: « Mathématiques », « Français »)
        \item \champ{Classe} — La classe concernée
        \item \champ{Enseignant} — L'enseignant responsable
        \item \champ{Volume horaire} — Nombre d'heures par semaine
        \item \champ{Coefficient} — Le coefficient de la matière
      \end{itemize}
    \item Cliquez sur \bouton{Enregistrer}
  \end{enumerate}
\end{etape}

\screenshot{admin_cours.png}{Gestion des cours et matières}

\subsection{Programme d'évaluation}

\begin{etape}{\faStepForward\quad Créer une évaluation}
  \begin{enumerate}[label=\textbf{\arabic*.}]
    \item Dans l'onglet \textbf{Évaluations}, cliquez sur \boutonvert{+ Nouvelle évaluation}
    \item Sélectionnez le \champ{Cours}, la \champ{Classe}, le \champ{Trimestre}
    \item Indiquez le \champ{Type} (Devoir, Examen, Interrogation)
    \item Indiquez la \champ{Date} et le \champ{Barème} (note sur 20, sur 10, etc.)
    \item Cliquez sur \bouton{Créer}
  \end{enumerate}
\end{etape}

\screenshot{admin_evaluations.png}{Création d'une évaluation}

% ─────────────────────────────────────────────────────────────────────────────
\section{Emploi du Temps}
\label{sec:admin_emploi}

\begin{etape}{\faStepForward\quad Configurer l'emploi du temps}
  \begin{enumerate}[label=\textbf{\arabic*.}]
    \item Dans le menu, cliquez sur \menu{Emploi du temps}
    \item La grille horaire s'affiche avec les jours en colonnes et les heures en lignes
    \item Pour ajouter un créneau :
      \begin{itemize}
        \item Cliquez sur la case vide correspondant au jour et à l'heure souhaités
        \item Sélectionnez le \champ{Cours}, la \champ{Classe} et la \champ{Salle}
        \item Cliquez sur \bouton{Ajouter}
      \end{itemize}
    \item Le créneau apparaît dans la grille avec un code couleur
  \end{enumerate}
\end{etape}

\screenshot{admin_emploi_temps.png}{Emploi du temps — Grille hebdomadaire}

\begin{astuce}
  Utilisez le filtre \textbf{Classe} en haut de la page pour afficher l'emploi du temps d'une classe spécifique.
\end{astuce}

% ─────────────────────────────────────────────────────────────────────────────
\section{Gestion des Parents}
\label{sec:admin_parents}

\begin{etape}{\faStepForward\quad Ajouter un parent}
  \begin{enumerate}[label=\textbf{\arabic*.}]
    \item Dans le menu, cliquez sur \menu{Parents}
    \item Cliquez sur \boutonvert{+ Ajouter un parent}
    \item Remplissez le formulaire :
      \begin{itemize}
        \item \champ{Nom} et \champ{Prénom} du parent
        \item \champ{Téléphone} et \champ{Email}
        \item \champ{Profession}
        \item \champ{Adresse}
        \item \champ{Enfant(s)} — Sélectionnez le ou les élèves liés à ce parent
      \end{itemize}
    \item Cliquez sur \bouton{Enregistrer}
    \item Le parent recevra ses identifiants de connexion
  \end{enumerate}
\end{etape}

\screenshot{admin_parents.png}{Gestion des parents — Liste et ajout}

% ─────────────────────────────────────────────────────────────────────────────
\section{Paiements}
\label{sec:admin_paiements}

\begin{etape}{\faStepForward\quad Enregistrer un paiement}
  \begin{enumerate}[label=\textbf{\arabic*.}]
    \item Dans le menu, cliquez sur \menu{Paiements}
    \item La page affiche les statistiques de paiement et la liste des paiements effectués
    \item Cliquez sur \boutonvert{+ Nouveau paiement}
    \item Remplissez :
      \begin{itemize}
        \item \champ{Élève} — Sélectionnez l'élève dans la liste
        \item \champ{Tranche} — Sélectionnez la tranche de scolarité (1ère, 2ème, 3ème tranche)
        \item \champ{Montant} — Le montant payé en FCFA
        \item \champ{Mode de paiement} — Espèces, Virement, Mobile Money
      \end{itemize}
    \item Cliquez sur \bouton{Enregistrer le paiement}
    \item Un reçu de paiement est automatiquement généré
  \end{enumerate}
\end{etape}

\screenshot{admin_paiements.png}{Module de paiements — Statistiques et enregistrement}

\screenshot{admin_paiements_recu.png}{Reçu de paiement généré automatiquement}

% ─────────────────────────────────────────────────────────────────────────────
\section{Bus Scolaire}
\label{sec:admin_bus}

\begin{etape}{\faStepForward\quad Gérer le transport scolaire}
  \begin{enumerate}[label=\textbf{\arabic*.}]
    \item Dans le menu, cliquez sur \menu{Bus Scolaire}
    \item La page affiche la liste des bus, leurs itinéraires et les élèves inscrits
    \item Pour ajouter un bus, cliquez sur \boutonvert{+ Nouveau bus}
    \item Renseignez le \champ{Numéro d'immatriculation}, la \champ{Capacité} et l'\champ{Itinéraire}
    \item Pour inscrire un élève au bus, cliquez sur \bouton{Ajouter un passager}
  \end{enumerate}
\end{etape}

\screenshot{admin_bus.png}{Gestion du transport scolaire}

% ─────────────────────────────────────────────────────────────────────────────
\section{Bulletins}
\label{sec:admin_bulletins}

\begin{etape}{\faStepForward\quad Générer et publier les bulletins}
  \begin{enumerate}[label=\textbf{\arabic*.}]
    \item Dans le menu, cliquez sur \menu{Bulletins}
    \item Sélectionnez la \champ{Classe}, le \champ{Trimestre} et l'\champ{Année académique}
    \item Cliquez sur \bouton{Générer les bulletins}
    \item Le système calcule automatiquement les moyennes, rangs et appréciations
    \item Vérifiez les bulletins générés dans le tableau
    \item Cliquez sur \boutonvert{Publier} pour rendre les bulletins visibles aux parents
  \end{enumerate}
\end{etape}

\screenshot{admin_bulletins_liste.png}{Liste des bulletins générés par classe}

\screenshot{admin_bulletin_detail.png}{Détail d'un bulletin — Notes, moyennes, rang et appréciations}

\begin{note}
  Un bulletin publié est visible par le parent dans son espace personnel. Avant publication, seuls les administrateurs et directeurs peuvent le consulter.
\end{note}

% ─────────────────────────────────────────────────────────────────────────────
\section{Discipline et Absences}
\label{sec:admin_discipline}

\begin{etape}{\faStepForward\quad Gérer la discipline}
  \begin{enumerate}[label=\textbf{\arabic*.}]
    \item Dans le menu, cliquez sur \menu{Discipline \& Absences}
    \item Le module comporte deux onglets :
      \begin{itemize}
        \item \textbf{Absences} — Liste des élèves absents avec la date et la durée
        \item \textbf{Sanctions} — Avertissements, blâmes et exclusions
      \end{itemize}
    \item Pour enregistrer une absence, cliquez sur \boutonvert{+ Signaler une absence}
    \item Sélectionnez l'\champ{Élève}, la \champ{Date} et le \champ{Motif}
    \item Cliquez sur \bouton{Enregistrer}
  \end{enumerate}
\end{etape}

\screenshot{admin_discipline.png}{Module Discipline \& Absences}

% ─────────────────────────────────────────────────────────────────────────────
\section{Personnel}
\label{sec:admin_personnel}

\begin{etape}{\faStepForward\quad Gérer le personnel}
  \begin{enumerate}[label=\textbf{\arabic*.}]
    \item Dans le menu, cliquez sur \menu{Personnel}
    \item La page affiche le personnel non-enseignant (agents de sécurité, secrétaires, etc.)
    \item Ajoutez, modifiez ou supprimez les membres du personnel
  \end{enumerate}
\end{etape}

\screenshot{admin_personnel.png}{Gestion du personnel administratif}

% ─────────────────────────────────────────────────────────────────────────────
\section{Rapports}
\label{sec:admin_rapports}

\begin{etape}{\faStepForward\quad Consulter les rapports}
  \begin{enumerate}[label=\textbf{\arabic*.}]
    \item Dans le menu, cliquez sur \menu{Rapports}
    \item Sélectionnez le type de rapport : \textbf{Pédagogique}, \textbf{Financier} ou \textbf{Discipline}
    \item Les rapports sont générés automatiquement à partir des données de l'application
    \item Vous pouvez les \textbf{imprimer} ou les \textbf{exporter en PDF}
  \end{enumerate}
\end{etape}

\screenshot{admin_rapports.png}{Module des rapports — Sélection et affichage}

% ─────────────────────────────────────────────────────────────────────────────
\section{Messagerie}
\label{sec:admin_messagerie}

\begin{etape}{\faStepForward\quad Envoyer un message}
  \begin{enumerate}[label=\textbf{\arabic*.}]
    \item Dans le menu, cliquez sur \menu{Messages}
    \item La messagerie s'ouvre avec :
      \begin{itemize}
        \item À gauche : la \textbf{liste des contacts} (enseignants, parents, directeurs, intendants)
        \item À droite : la \textbf{conversation} sélectionnée
      \end{itemize}
    \item Cliquez sur un contact pour ouvrir la conversation
    \item Tapez votre message dans la zone de texte en bas
    \item Cliquez sur \bouton{Envoyer} ou appuyez sur Entrée
  \end{enumerate}
\end{etape}

\screenshot{admin_messagerie.png}{Interface de messagerie — Conversations et envoi}

\begin{note}
  Le badge rouge \textcolor{danger}{\faBell} à côté de « Messages » dans le menu indique le nombre de messages non lus. Ce compteur se met à jour automatiquement.
\end{note}

% ─────────────────────────────────────────────────────────────────────────────
\section{Mon Profil}
\label{sec:admin_profil}

\begin{etape}{\faStepForward\quad Consulter/modifier son profil}
  \begin{enumerate}[label=\textbf{\arabic*.}]
    \item Dans le menu, cliquez sur \menu{Mon Profil}
    \item Votre fiche personnelle s'affiche avec vos informations
    \item Vous pouvez modifier votre \champ{Mot de passe}, votre \champ{Email} et votre \champ{Téléphone}
    \item Cliquez sur \bouton{Enregistrer} pour sauvegarder les modifications
  \end{enumerate}
\end{etape}

\screenshot{admin_profil.png}{Page de profil administrateur}


% ============================================================================
%  CHAPITRE 4 : ESPACE DIRECTEUR
% ============================================================================
\chapter{Espace Directeur}
\label{chap:directeur}

Le directeur a une vue de \textbf{supervision} sur l'ensemble de l'établissement. Il peut consulter toutes les données académiques et financières sans nécessairement les modifier.

\section{Tableau de bord Directeur}

\begin{etape}{\faStepForward\quad Accès}
  Après connexion avec un compte \role{Directeur}, vous arrivez automatiquement sur le tableau de bord directeur.
\end{etape}

\screenshot{directeur_dashboard.png}{Tableau de bord Directeur — Supervision globale}

\subsection*{Fonctionnalités du tableau de bord}
\begin{itemize}[label=\textcolor{brandOrange}{\faCheckCircle}]
  \item Synthèse des effectifs (élèves, enseignants, classes)
  \item État des paiements (montant collecté, taux de recouvrement)
  \item Dernières inscriptions
  \item Alertes et notifications
\end{itemize}

\section{Menu du Directeur}

Le directeur dispose des sections suivantes dans son menu :

\begin{table}[htbp]
  \centering
  \caption{Sections accessibles au Directeur}
  \begin{tabularx}{\textwidth}{lX}
    \toprule
    \textbf{Section} & \textbf{Description} \\
    \midrule
    Tableau de bord & Vue d'ensemble avec statistiques \\
    Messages & Messagerie interne — communication avec tous les acteurs \\
    Rapports & Consultation des rapports pédagogiques et financiers \\
    Élèves & Consultation de la liste des élèves (lecture seule) \\
    Enseignants & Consultation de la liste des enseignants \\
    Classes & Consultation des classes \\
    Salles & Consultation des salles \\
    Paiements & Suivi des paiements \\
    Tranches & Gestion des tranches de scolarité \\
    Élèves en impayé & Liste des élèves n'ayant pas soldé leur scolarité \\
    Mon Profil & Consultation et modification du profil personnel \\
    \bottomrule
  \end{tabularx}
\end{table}

\section{Suivi des Paiements}

\begin{etape}{\faStepForward\quad Consulter les paiements}
  \begin{enumerate}[label=\textbf{\arabic*.}]
    \item Cliquez sur \menu{Paiements} dans le menu
    \item La page affiche :
      \begin{itemize}
        \item Le montant total collecté
        \item Le nombre de paiements par tranche
        \item La liste détaillée des paiements
      \end{itemize}
    \item Utilisez les filtres pour cibler une classe ou une période
  \end{enumerate}
\end{etape}

\screenshot{directeur_paiements.png}{Vue des paiements — Espace Directeur}

\section{Élèves en impayé}

\begin{etape}{\faStepForward\quad Identifier les impayés}
  \begin{enumerate}[label=\textbf{\arabic*.}]
    \item Cliquez sur \menu{Élèves en impayé}
    \item La liste affiche les élèves dont la scolarité n'est pas entièrement soldée
    \item Pour chaque élève : le montant payé, le montant restant et le pourcentage de paiement
  \end{enumerate}
\end{etape}

\screenshot{directeur_impayes.png}{Liste des élèves en impayé}

\section{Communication}

Le directeur peut communiquer avec \textbf{tous les acteurs} de l'application via la messagerie interne. La procédure est identique à celle décrite dans la section~\ref{sec:admin_messagerie}.

\screenshot{directeur_messages.png}{Messagerie du Directeur}


% ============================================================================
%  CHAPITRE 5 : ESPACE INTENDANT
% ============================================================================
\chapter{Espace Intendant}
\label{chap:intendant}

L'intendant est responsable de la \textbf{gestion financière} de l'établissement. Son espace est centré sur les paiements, les tranches de scolarité et le suivi des impayés.

\section{Tableau de bord Intendant}

\screenshot{intendant_dashboard.png}{Tableau de bord Intendant — Vue financière}

\subsection*{Indicateurs financiers}
\begin{itemize}[label=\textcolor{brandOrange}{\faMoneyBillWave}]
  \item Montant total collecté
  \item Nombre de paiements du jour
  \item Taux de recouvrement global
  \item Montant des impayés
\end{itemize}

\section{Enregistrer un paiement}

La procédure est identique à celle décrite dans la section~\ref{sec:admin_paiements}. L'intendant accède à la même interface de paiements.

\screenshot{intendant_paiement.png}{Enregistrement d'un paiement — Espace Intendant}

\section{Gérer les tranches de scolarité}

\begin{etape}{\faStepForward\quad Configurer les tranches}
  \begin{enumerate}[label=\textbf{\arabic*.}]
    \item Cliquez sur \menu{Tranches} dans le menu
    \item La page affiche les tranches existantes avec les montants et les dates limites
    \item Pour créer une nouvelle tranche, cliquez sur \boutonvert{+ Nouvelle tranche}
    \item Renseignez :
      \begin{itemize}
        \item \champ{Libellé} — Ex : « 1ère Tranche », « 2ème Tranche »
        \item \champ{Montant} — Le montant en FCFA
        \item \champ{Date limite} — La date limite de paiement
        \item \champ{Classe} — La classe concernée (ou toutes les classes)
      \end{itemize}
    \item Cliquez sur \bouton{Enregistrer}
  \end{enumerate}
\end{etape}

\screenshot{intendant_tranches.png}{Gestion des tranches de scolarité}

\section{Suivi des impayés}

\begin{etape}{\faStepForward\quad Consulter les impayés}
  \begin{enumerate}[label=\textbf{\arabic*.}]
    \item Cliquez sur \menu{Élèves en impayé}
    \item La liste affiche les élèves avec le solde restant
    \item Vous pouvez filtrer par classe ou par tranche
    \item Utilisez le bouton \bouton{Exporter} pour générer un fichier des impayés
  \end{enumerate}
\end{etape}

\screenshot{intendant_impayes.png}{Suivi des impayés — Espace Intendant}


% ============================================================================
%  CHAPITRE 6 : ESPACE ENSEIGNANT
% ============================================================================
\chapter{Espace Enseignant}
\label{chap:enseignant}

L'enseignant dispose d'un espace personnel dédié à ses activités pédagogiques.

\section{Tableau de bord Enseignant}

\begin{etape}{\faStepForward\quad Accès}
  Après connexion avec un compte \role{Enseignant}, vous arrivez automatiquement sur votre tableau de bord personnel.
\end{etape}

\screenshot{enseignant_dashboard.png}{Tableau de bord Enseignant — Vue personnelle}

\subsection*{Informations affichées}
\begin{itemize}[label=\textcolor{success}{\faCheckCircle}]
  \item Nombre de classes assignées
  \item Nombre total d'élèves
  \item Prochains cours de la journée
  \item Statistiques d'évaluation
\end{itemize}

\section{Menu de l'enseignant}

\begin{table}[htbp]
  \centering
  \caption{Sections accessibles à l'Enseignant}
  \begin{tabularx}{\textwidth}{lX}
    \toprule
    \textbf{Section} & \textbf{Description} \\
    \midrule
    Dashboard Enseignant & Vue d'ensemble personnelle \\
    Messages & Messagerie — communiquer avec l'administration, les parents \\
    Emploi du temps & Consulter son emploi du temps hebdomadaire \\
    Notes \& Évaluations & Saisir les notes et créer des évaluations \\
    Paramètres & Modifier son profil et mot de passe \\
    \bottomrule
  \end{tabularx}
\end{table}

\section{Consulter son emploi du temps}

\begin{etape}{\faStepForward\quad Voir l'emploi du temps}
  \begin{enumerate}[label=\textbf{\arabic*.}]
    \item Cliquez sur \menu{Emploi du temps} dans le menu
    \item Votre emploi du temps hebdomadaire s'affiche sous forme de grille
    \item Chaque créneau indique : la matière, la classe et la salle
    \item Les créneaux sont colorés pour faciliter la lecture
  \end{enumerate}
\end{etape}

\screenshot{enseignant_emploi.png}{Emploi du temps personnel de l'enseignant}

\section{Saisir les notes}

\begin{etape}{\faStepForward\quad Saisie des notes}
  \begin{enumerate}[label=\textbf{\arabic*.}]
    \item Cliquez sur \menu{Notes \& Évaluations}
    \item Sélectionnez la \champ{Classe} dans la liste déroulante
    \item Sélectionnez l'\champ{Évaluation} (Devoir, Examen, etc.)
    \item La liste des élèves de la classe s'affiche avec des champs de saisie
    \item Pour chaque élève, saisissez la note dans le champ correspondant
    \item Cliquez sur \boutonvert{Enregistrer les notes}
  \end{enumerate}
\end{etape}

\screenshot{enseignant_saisie_notes.png}{Interface de saisie des notes par l'enseignant}

\begin{attention}
  Vérifiez bien les notes avant de les enregistrer. Une fois enregistrées, seul un administrateur peut les modifier.
\end{attention}

\section{Créer une évaluation}

\begin{etape}{\faStepForward\quad Créer une nouvelle évaluation}
  \begin{enumerate}[label=\textbf{\arabic*.}]
    \item Dans la page \textbf{Notes \& Évaluations}, cliquez sur \boutonvert{+ Nouvelle évaluation}
    \item Remplissez :
      \begin{itemize}
        \item \champ{Titre} — Ex : « Devoir de Mathématiques n°3 »
        \item \champ{Classe} — La classe concernée
        \item \champ{Trimestre} — Le trimestre en cours
        \item \champ{Type} — Devoir, Examen, Interrogation
        \item \champ{Date} — La date de l'évaluation
        \item \champ{Barème} — Note maximale (ex : 20)
      \end{itemize}
    \item Cliquez sur \bouton{Créer l'évaluation}
  \end{enumerate}
\end{etape}

\screenshot{enseignant_creation_eval.png}{Formulaire de création d'évaluation}

\section{Messagerie de l'enseignant}

\begin{etape}{\faStepForward\quad Communiquer}
  \begin{enumerate}[label=\textbf{\arabic*.}]
    \item Cliquez sur \menu{Messages} dans le menu
    \item Vous pouvez envoyer des messages à :
      \begin{itemize}
        \item L'administrateur
        \item Le directeur
        \item Les parents de vos élèves
        \item Les autres enseignants
      \end{itemize}
    \item Sélectionnez un contact, rédigez votre message et cliquez sur \bouton{Envoyer}
  \end{enumerate}
\end{etape}

\screenshot{enseignant_messagerie.png}{Messagerie de l'enseignant}

\begin{astuce}
  Le badge rouge sur l'icône \textcolor{danger}{\faEnvelope} « Messages » vous indique combien de messages non lus vous avez. Les notifications se mettent à jour automatiquement.
\end{astuce}


% ============================================================================
%  CHAPITRE 7 : ESPACE PARENT
% ============================================================================
\chapter{Espace Parent}
\label{chap:parent}

L'espace parent permet aux parents d'élèves de \textbf{suivre la scolarité de leur enfant} à distance.

\section{Tableau de bord Parent}

\begin{etape}{\faStepForward\quad Accès}
  Après connexion avec un compte \role{Parent}, vous arrivez sur le tableau de bord parent.
\end{etape}

\screenshot{parent_dashboard.png}{Tableau de bord Parent — Suivi de l'enfant}

\subsection*{Informations affichées}
\begin{itemize}[label=\textcolor{brandBlue}{\faChild}]
  \item Nom et photo de l'enfant
  \item Classe et salle actuelles
  \item Dernières notes et moyenne
  \item État des paiements
  \item Absences récentes
\end{itemize}

\section{Menu du parent}

\begin{table}[htbp]
  \centering
  \caption{Sections accessibles au Parent}
  \begin{tabularx}{\textwidth}{lX}
    \toprule
    \textbf{Section} & \textbf{Description} \\
    \midrule
    Dashboard Parent & Vue d'ensemble sur l'enfant \\
    Messages & Communiquer avec l'école (enseignants, administration) \\
    Paiements & Consulter l'état des paiements de scolarité \\
    Rapports & Consulter les rapports et bulletins \\
    Paramètres & Modifier son profil et mot de passe \\
    \bottomrule
  \end{tabularx}
\end{table}

\section{Consulter les résultats scolaires}

\begin{etape}{\faStepForward\quad Voir les notes}
  \begin{enumerate}[label=\textbf{\arabic*.}]
    \item Sur le tableau de bord, la section « Dernières notes » affiche les notes récentes
    \item Cliquez sur \bouton{Voir le bulletin} pour accéder au bulletin complet
    \item Le bulletin affiche :
      \begin{itemize}
        \item Toutes les matières avec les notes et coefficients
        \item La moyenne générale
        \item Le rang dans la classe
        \item L'appréciation du conseil de classe
      \end{itemize}
  \end{enumerate}
\end{etape}

\screenshot{parent_bulletin.png}{Bulletin scolaire vu par le parent}

\section{Consulter les paiements}

\begin{etape}{\faStepForward\quad Voir l'état des paiements}
  \begin{enumerate}[label=\textbf{\arabic*.}]
    \item Cliquez sur \menu{Paiements} dans le menu
    \item La page affiche :
      \begin{itemize}
        \item Le montant total de la scolarité
        \item Le montant déjà payé
        \item Le montant restant à payer
        \item L'historique des paiements avec les dates et reçus
      \end{itemize}
  \end{enumerate}
\end{etape}

\screenshot{parent_paiements.png}{État des paiements — Espace Parent}

\section{Communiquer avec l'école}

\begin{etape}{\faStepForward\quad Envoyer un message}
  \begin{enumerate}[label=\textbf{\arabic*.}]
    \item Cliquez sur \menu{Messages}
    \item Sélectionnez le destinataire (enseignant, directeur, administration)
    \item Rédigez votre message
    \item Cliquez sur \bouton{Envoyer}
  \end{enumerate}
\end{etape}

\screenshot{parent_messagerie.png}{Messagerie — Espace Parent}

\begin{astuce}
  Utilisez la messagerie pour :
  \begin{itemize}
    \item Justifier une absence de votre enfant
    \item Demander un rendez-vous avec un enseignant
    \item Poser une question sur les résultats scolaires
    \item Signaler un problème
  \end{itemize}
\end{astuce}


% ============================================================================
%  CHAPITRE 8 : CAS D'UTILISATION DÉTAILLÉS
% ============================================================================
\chapter{Cas d'Utilisation Détaillés}
\label{chap:cas_utilisation}

Ce chapitre présente les \textbf{scénarios d'utilisation courants} sous forme de tableaux récapitulatifs, afin de servir de référence rapide.

\section{Cas d'utilisation — Administrateur}

\begin{table}[htbp]
  \centering
  \caption{Cas d'utilisation de l'Administrateur}
  \small
  \rowcolors{2}{bgLight}{white}
  \begin{tabularx}{\textwidth}{>{\bfseries}p{3.5cm} p{3cm} X p{2.5cm}}
    \toprule
    \textbf{Cas d'utilisation} & \textbf{Pré-condition} & \textbf{Description} & \textbf{Résultat} \\
    \midrule
    Créer un élève & Être connecté comme admin & Saisir les informations personnelles de l'élève et l'affecter à une classe & Élève créé avec un matricule unique \\
    Inscrire un élève & Élève existant, classe existante & Associer un élève à une classe et une salle pour l'année en cours & Inscription enregistrée \\
    Créer un enseignant & Être connecté comme admin & Saisir les informations et le cours de l'enseignant & Compte enseignant créé avec identifiants \\
    Affecter un enseignant & Enseignant et classe existants & Lier un enseignant à une ou plusieurs classes & Emploi du temps mis à jour \\
    Créer une classe & Être connecté comme admin & Définir un libellé et un cycle & Classe disponible dans le système \\
    Créer une salle & Classe existante & Définir une salle avec sa capacité & Salle disponible pour les inscriptions \\
    Créer un cours & Classe et enseignant existants & Définir une matière avec coefficient et volume horaire & Cours créé et assigné \\
    Planifier l'emploi du temps & Cours et salles existants & Placer les créneaux horaires dans la grille & Emploi du temps publié \\
    Enregistrer un paiement & Élève inscrit & Saisir le montant et le mode de paiement & Reçu de paiement généré \\
    Générer un bulletin & Notes saisies, trimestre configuré & Calculer les moyennes et rangs automatiquement & Bulletin prêt à publier \\
    Publier un bulletin & Bulletin généré & Rendre le bulletin visible aux parents & Parents notifiés \\
    Enregistrer une absence & Élève existant & Saisir la date, l'heure et le motif & Absence enregistrée \\
    Envoyer un message & Être connecté & Rédiger et envoyer un message à un autre acteur & Message délivré avec notification \\
    Gérer le bus scolaire & Être connecté comme admin & Créer un bus et y inscrire des élèves & Élèves affectés au bus \\
    \bottomrule
  \end{tabularx}
\end{table}

\section{Cas d'utilisation — Directeur}

\begin{table}[htbp]
  \centering
  \caption{Cas d'utilisation du Directeur}
  \small
  \rowcolors{2}{bgLight}{white}
  \begin{tabularx}{\textwidth}{>{\bfseries}p{3.5cm} p{3cm} X p{2.5cm}}
    \toprule
    \textbf{Cas d'utilisation} & \textbf{Pré-condition} & \textbf{Description} & \textbf{Résultat} \\
    \midrule
    Consulter le tableau de bord & Être connecté comme directeur & Voir les statistiques globales de l'établissement & Vue d'ensemble affichée \\
    Consulter les élèves & Être connecté & Parcourir la liste des élèves par classe & Liste affichée \\
    Consulter les enseignants & Être connecté & Voir les enseignants actifs et leurs matières & Liste affichée \\
    Suivre les paiements & Être connecté & Consulter les montants collectés et les impayés & Données financières affichées \\
    Identifier les impayés & Être connecté & Voir la liste des élèves en retard de paiement & Liste des impayés \\
    Consulter les rapports & Être connecté & Lire les rapports pédagogiques et financiers & Rapport affiché \\
    Envoyer un message & Être connecté & Communiquer avec enseignants, parents, admin, intendant & Message délivré \\
    Modifier son profil & Être connecté & Changer son mot de passe ou ses informations & Profil mis à jour \\
    \bottomrule
  \end{tabularx}
\end{table}

\section{Cas d'utilisation — Intendant}

\begin{table}[htbp]
  \centering
  \caption{Cas d'utilisation de l'Intendant}
  \small
  \rowcolors{2}{bgLight}{white}
  \begin{tabularx}{\textwidth}{>{\bfseries}p{3.5cm} p{3cm} X p{2.5cm}}
    \toprule
    \textbf{Cas d'utilisation} & \textbf{Pré-condition} & \textbf{Description} & \textbf{Résultat} \\
    \midrule
    Consulter le tableau de bord & Être connecté comme intendant & Voir les indicateurs financiers & Vue financière affichée \\
    Enregistrer un paiement & Élève inscrit & Saisir le montant, la tranche et le mode & Paiement enregistré et reçu généré \\
    Configurer les tranches & Être connecté & Définir les montants et dates limites par tranche & Tranches configurées \\
    Suivre les impayés & Être connecté & Identifier les élèves en retard de paiement & Liste des impayés \\
    Consulter les inscriptions & Être connecté & Vérifier les inscriptions et les scolarités & Liste affichée \\
    Envoyer un message & Être connecté & Communiquer avec les autres acteurs & Message délivré \\
    \bottomrule
  \end{tabularx}
\end{table}

\section{Cas d'utilisation — Enseignant}

\begin{table}[htbp]
  \centering
  \caption{Cas d'utilisation de l'Enseignant}
  \small
  \rowcolors{2}{bgLight}{white}
  \begin{tabularx}{\textwidth}{>{\bfseries}p{3.5cm} p{3cm} X p{2.5cm}}
    \toprule
    \textbf{Cas d'utilisation} & \textbf{Pré-condition} & \textbf{Description} & \textbf{Résultat} \\
    \midrule
    Consulter le tableau de bord & Être connecté comme enseignant & Voir ses classes, élèves, prochains cours & Vue personnelle affichée \\
    Consulter l'emploi du temps & Être connecté & Voir sa grille horaire hebdomadaire & Emploi du temps affiché \\
    Saisir les notes & Évaluation créée & Entrer les notes de chaque élève pour une évaluation & Notes enregistrées \\
    Créer une évaluation & Être connecté, cours existant & Définir le titre, type, date et barème & Évaluation créée \\
    Faire l'appel & Être connecté, créneau de cours & Marquer les élèves présents et absents & Présences enregistrées \\
    Envoyer un message & Être connecté & Communiquer avec l'administration, les parents & Message délivré \\
    Modifier son profil & Être connecté & Changer son mot de passe & Profil mis à jour \\
    \bottomrule
  \end{tabularx}
\end{table}

\section{Cas d'utilisation — Parent}

\begin{table}[htbp]
  \centering
  \caption{Cas d'utilisation du Parent}
  \small
  \rowcolors{2}{bgLight}{white}
  \begin{tabularx}{\textwidth}{>{\bfseries}p{3.5cm} p{3cm} X p{2.5cm}}
    \toprule
    \textbf{Cas d'utilisation} & \textbf{Pré-condition} & \textbf{Description} & \textbf{Résultat} \\
    \midrule
    Consulter le tableau de bord & Être connecté comme parent & Voir les infos de l'enfant, notes récentes, absences & Vue enfant affichée \\
    Consulter le bulletin & Bulletin publié & Voir les notes, la moyenne, le rang et l'appréciation & Bulletin affiché \\
    Consulter les paiements & Être connecté & Voir le montant payé, le restant et les reçus & État financier affiché \\
    Envoyer un message & Être connecté & Communiquer avec un enseignant ou l'administration & Message délivré \\
    Justifier une absence & Être connecté & Envoyer un message expliquant l'absence de l'enfant & Message envoyé \\
    Modifier son profil & Être connecté & Changer son mot de passe ou ses informations & Profil mis à jour \\
    \bottomrule
  \end{tabularx}
\end{table}


% ============================================================================
%  CHAPITRE 9 : FOIRE AUX QUESTIONS
% ============================================================================
\chapter{Foire Aux Questions (FAQ)}
\label{chap:faq}

\begin{tcolorbox}[colback=bgLight, colframe=brandBlue, rounded corners, title={\faQuestionCircle\quad Questions fréquentes}]

\textbf{Q1 : J'ai oublié mon mot de passe, que faire~?}\\
\textit{R :} Contactez l'administrateur de l'école. Lui seul peut réinitialiser votre mot de passe. Vous recevrez un nouveau mot de passe temporaire que vous devrez changer lors de votre prochaine connexion.\\[8pt]

\textbf{Q2 : La page ne s'affiche pas correctement, que faire~?}\\
\textit{R :} 
\begin{enumerate}[label=\arabic*.]
  \item Vérifiez votre connexion internet
  \item Essayez de rafraîchir la page (touche \textbf{F5} ou \textbf{Ctrl+R})
  \item Essayez un autre navigateur (Google Chrome est recommandé)
  \item Videz le cache de votre navigateur (\textbf{Ctrl+Maj+Suppr})
\end{enumerate}
\vspace{4pt}

\textbf{Q3 : Je suis parent et je ne vois pas le bulletin de mon enfant}\\
\textit{R :} Le bulletin doit d'abord être \textbf{publié} par l'administrateur ou le directeur. Si le trimestre est terminé mais que le bulletin n'apparaît pas, contactez l'école.\\[8pt]

\textbf{Q4 : Comment changer mon mot de passe~?}\\
\textit{R :} Allez dans \menu{Mon Profil} ou \menu{Paramètres} dans votre menu. Vous y trouverez l'option pour modifier votre mot de passe. Entrez l'ancien mot de passe, puis le nouveau (deux fois).\\[8pt]

\textbf{Q5 : L'application est très lente, que faire~?}\\
\textit{R :} 
\begin{itemize}
  \item Vérifiez votre connexion internet
  \item Fermez les autres onglets de votre navigateur
  \item Si le problème persiste, contactez le support technique
\end{itemize}
\vspace{4pt}

\textbf{Q6 : Je suis enseignant, comment accéder à mes cours~?}\\
\textit{R :} Après connexion, cliquez sur \menu{Emploi du temps} pour voir vos cours planifiés. Pour saisir les notes, allez dans \menu{Notes \& Évaluations}.\\[8pt]

\textbf{Q7 : Puis-je utiliser l'application sur mon téléphone~?}\\
\textit{R :} Oui ! L'application est \textbf{responsive} et s'adapte à la taille de votre écran. Ouvrez simplement votre navigateur mobile et tapez l'adresse de l'application.\\[8pt]

\textbf{Q8 : Les messages sont-ils confidentiels~?}\\
\textit{R :} Oui. Seuls l'expéditeur et le destinataire peuvent voir le contenu d'un message. L'administrateur système n'a pas accès aux messages privés entre les utilisateurs.

\end{tcolorbox}


% ============================================================================
%  CHAPITRE 10 : GLOSSAIRE
% ============================================================================
\chapter{Glossaire}
\label{chap:glossaire}

\begin{description}[style=nextline, leftmargin=3cm, font=\bfseries\color{brandBlue}]
  \item[Administrateur] Utilisateur ayant tous les droits sur l'application. Responsable de la configuration et de la gestion complète de l'école.
  \item[Année académique] Période scolaire allant généralement de septembre à juillet.
  \item[Bulletin] Document récapitulatif des notes et appréciations d'un élève pour un trimestre donné.
  \item[Classe] Regroupement d'élèves d'un même niveau (ex : SIL, CP, CE1, CE2, CM1, CM2).
  \item[Coefficient] Valeur numérique attribuée à une matière, reflétant son importance relative dans le calcul de la moyenne.
  \item[Cours] Matière enseignée (ex : Mathématiques, Français, Anglais).
  \item[Cycle] Regroupement de classes par niveau d'enseignement (ex : Cycle 1 = Maternel, Cycle 2 = Primaire).
  \item[Dashboard] Tableau de bord — page d'accueil affichant les indicateurs clés.
  \item[Directeur] Responsable pédagogique et administratif de l'établissement.
  \item[Emploi du temps] Planning hebdomadaire des cours pour une classe ou un enseignant.
  \item[Évaluation] Tout exercice noté (devoir, examen, interrogation) servant à évaluer les compétences d'un élève.
  \item[FCFA] Franc de la Communauté Financière Africaine — monnaie utilisée au Cameroun.
  \item[Fréquentation] Table du système qui enregistre la présence des élèves dans les salles.
  \item[Impayé] Montant de scolarité non encore réglé par un élève ou son parent.
  \item[Inscription] Acte administratif officialisant l'admission d'un élève dans une classe.
  \item[Intendant] Responsable financier de l'établissement.
  \item[Matricule] Identifiant unique attribué à chaque élève lors de son inscription.
  \item[Mobile Money] Service de paiement mobile (Orange Money, MTN Money).
  \item[Notification] Alerte visuelle (badge rouge) indiquant un événement nécessitant l'attention de l'utilisateur.
  \item[Parent] Tuteur légal d'un ou plusieurs élèves.
  \item[Salle] Espace physique où se déroulent les cours. Chaque salle est associée à une classe.
  \item[Sidebar] Menu de navigation latéral (à gauche de l'écran).
  \item[Tranche] Division du montant de la scolarité en plusieurs versements échelonnés.
  \item[Trimestre] Division de l'année scolaire en trois périodes (1er, 2e, 3e trimestre).
\end{description}


% ============================================================================
%  CHAPITRE 11 : SUPPORT TECHNIQUE
% ============================================================================
\chapter{Support Technique}
\label{chap:support}

\section{Contacter le support}

En cas de problème technique, vous pouvez contacter le support de plusieurs manières :

\begin{table}[htbp]
  \centering
  \caption{Moyens de contact du support technique}
  \begin{tabularx}{\textwidth}{lX}
    \toprule
    \textbf{Canal} & \textbf{Coordonnées} \\
    \midrule
    \faEnvelope\quad Email & support@ecolegest.cm \\
    \faPhone\quad Téléphone & +237 6XX XXX XXX \\
    \faWhatsapp\quad WhatsApp & +237 6XX XXX XXX \\
    \faMapMarkerAlt\quad Adresse & Melen, Yaoundé — Cameroun \\
    \faClock\quad Horaires & Lundi - Vendredi, 8h00 - 17h00 \\
    \bottomrule
  \end{tabularx}
\end{table}

\section{Informations à fournir au support}

Lorsque vous contactez le support technique, veuillez fournir :

\begin{enumerate}[label=\textbf{\arabic*.}]
  \item Votre \textbf{nom d'utilisateur}
  \item Votre \textbf{rôle} (admin, directeur, enseignant, parent, intendant)
  \item Une \textbf{description précise du problème}
  \item Une \textbf{capture d'écran} si possible (appuyez sur \textbf{Impr. écran} ou \textbf{Windows + Maj + S})
  \item Le \textbf{navigateur} que vous utilisez (Chrome, Firefox, etc.)
  \item L'\textbf{heure} à laquelle le problème s'est produit
\end{enumerate}

\section{Procédures de dépannage rapide}

\begin{table}[htbp]
  \centering
  \caption{Problèmes courants et solutions}
  \small
  \rowcolors{2}{bgLight}{white}
  \begin{tabularx}{\textwidth}{p{4cm} X}
    \toprule
    \textbf{Problème} & \textbf{Solution} \\
    \midrule
    Page blanche & Rafraîchissez la page (F5). Si le problème persiste, videz le cache du navigateur. \\
    Mot de passe oublié & Contactez l'administrateur pour une réinitialisation. \\
    Impossible de se connecter & Vérifiez que le nom d'utilisateur et le mot de passe sont corrects. Attention aux majuscules. \\
    Message non envoyé & Vérifiez votre connexion internet. Réessayez après quelques secondes. \\
    Bulletin non visible & Le bulletin n'a pas encore été publié. Contactez l'administration. \\
    Données non à jour & Rafraîchissez la page. Les données sont mises à jour en temps réel. \\
    Erreur 500 & Erreur serveur. Réessayez dans quelques minutes. Si le problème persiste, contactez le support. \\
    \bottomrule
  \end{tabularx}
\end{table}


% ============================================================================
%  ANNEXES
% ============================================================================
\appendix

\chapter{Raccourcis et Astuces}
\label{app:raccourcis}

\begin{table}[htbp]
  \centering
  \caption{Raccourcis clavier utiles}
  \begin{tabularx}{\textwidth}{lX}
    \toprule
    \textbf{Raccourci} & \textbf{Action} \\
    \midrule
    \textbf{F5} ou \textbf{Ctrl+R} & Rafraîchir la page \\
    \textbf{Ctrl+F} & Rechercher du texte dans la page \\
    \textbf{Ctrl+P} & Imprimer la page \\
    \textbf{Ctrl+Maj+Suppr} & Vider le cache du navigateur \\
    \textbf{Entrée} & Valider un formulaire ou envoyer un message \\
    \textbf{Tab} & Passer au champ suivant dans un formulaire \\
    \textbf{Échap} & Fermer une fenêtre modale ou un popup \\
    \bottomrule
  \end{tabularx}
\end{table}

\chapter{Index des Captures d'Écran}
\label{app:index_screenshots}

Le tableau ci-dessous récapitule toutes les captures d'écran nécessaires pour compléter ce manuel. Placez les fichiers images dans le dossier \texttt{screenshots/} à côté du fichier \texttt{.tex}.

\begin{table}[htbp]
  \centering
  \caption{Liste des captures d'écran à réaliser}
  \small
  \rowcolors{2}{bgLight}{white}
  \begin{tabularx}{\textwidth}{>{\ttfamily\small}p{5.5cm} X l}
    \toprule
    \textbf{Nom du fichier} & \textbf{Description} & \textbf{Page} \\
    \midrule
    logo.png & Logo de l'application (512×512 px) & Couverture \\
    schema\_general.png & Schéma des modules et acteurs & Chap. 1 \\
    landing\_page.png & Page d'accueil publique & Chap. 2 \\
    login\_page.png & Formulaire de connexion & Chap. 2 \\
    login\_success.png & Redirection après connexion & Chap. 2 \\
    sidebar\_deconnexion.png & Menu latéral avec bouton déconnexion & Chap. 2 \\
    interface\_zones.png & Zones de l'interface annotées & Chap. 2 \\
    admin\_dashboard.png & Tableau de bord admin & Chap. 3 \\
    admin\_eleves\_liste.png & Liste des élèves & Chap. 3 \\
    admin\_eleves\_ajout.png & Formulaire ajout élève & Chap. 3 \\
    admin\_enseignants\_liste.png & Liste des enseignants & Chap. 3 \\
    admin\_enseignants\_ajout.png & Formulaire ajout enseignant & Chap. 3 \\
    admin\_enseignants\_affectation.png & Affectation enseignant & Chap. 3 \\
    admin\_classes.png & Gestion des classes & Chap. 3 \\
    admin\_salles.png & Gestion des salles & Chap. 3 \\
    admin\_inscriptions.png & Module d'inscription & Chap. 3 \\
    admin\_cours.png & Gestion des cours & Chap. 3 \\
    admin\_evaluations.png & Création d'évaluation & Chap. 3 \\
    admin\_emploi\_temps.png & Emploi du temps & Chap. 3 \\
    admin\_parents.png & Gestion des parents & Chap. 3 \\
    admin\_paiements.png & Module de paiements & Chap. 3 \\
    admin\_paiements\_recu.png & Reçu de paiement & Chap. 3 \\
    admin\_bus.png & Gestion du bus & Chap. 3 \\
    admin\_bulletins\_liste.png & Liste des bulletins & Chap. 3 \\
    admin\_bulletin\_detail.png & Détail d'un bulletin & Chap. 3 \\
    admin\_discipline.png & Discipline et absences & Chap. 3 \\
    admin\_personnel.png & Gestion du personnel & Chap. 3 \\
    admin\_rapports.png & Module des rapports & Chap. 3 \\
    admin\_messagerie.png & Interface de messagerie & Chap. 3 \\
    admin\_profil.png & Profil administrateur & Chap. 3 \\
    directeur\_dashboard.png & Tableau de bord directeur & Chap. 4 \\
    directeur\_paiements.png & Paiements vu par le directeur & Chap. 4 \\
    directeur\_impayes.png & Impayés vu par le directeur & Chap. 4 \\
    directeur\_messages.png & Messagerie du directeur & Chap. 4 \\
    intendant\_dashboard.png & Tableau de bord intendant & Chap. 5 \\
    intendant\_paiement.png & Paiement vu par l'intendant & Chap. 5 \\
    intendant\_tranches.png & Gestion des tranches & Chap. 5 \\
    intendant\_impayes.png & Impayés vu par l'intendant & Chap. 5 \\
    enseignant\_dashboard.png & Tableau de bord enseignant & Chap. 6 \\
    enseignant\_emploi.png & Emploi du temps enseignant & Chap. 6 \\
    enseignant\_saisie\_notes.png & Saisie des notes & Chap. 6 \\
    enseignant\_creation\_eval.png & Création d'évaluation & Chap. 6 \\
    enseignant\_messagerie.png & Messagerie enseignant & Chap. 6 \\
    parent\_dashboard.png & Tableau de bord parent & Chap. 7 \\
    parent\_bulletin.png & Bulletin vu par le parent & Chap. 7 \\
    parent\_paiements.png & Paiements vu par le parent & Chap. 7 \\
    parent\_messagerie.png & Messagerie parent & Chap. 7 \\
    \bottomrule
  \end{tabularx}
\end{table}

% ============================================================================
%  FIN DU DOCUMENT
% ============================================================================

\vfill
\begin{center}
  \rule{6cm}{0.5pt}\\[10pt]
  {\large\textbf{EcoleGest — Application de Gestion Scolaire}}\\[4pt]
  {\small Melen, Yaoundé — Cameroun}\\[4pt]
  {\small\textcopyright{} 2026 — Tous droits réservés}
\end{center}

\end{document}
